% 第5章 現代ポートフォリオ理論(1)(資料5 + Markowitz補足)
\chapter{現代ポートフォリオ理論 (1)：分散投資と平均分散モデル}

本章から3章にわたり，\textbf{現代ポートフォリオ理論} (Modern Portfolio Theory)
を学ぶ。本章では，ポートフォリオの収益率の定式化，分散投資によるリスク低減効果
(分散投資効果)，そしてリスク性資産のみからなる場合の平均分散モデル
(Markowitzのポートフォリオ選択問題)を扱う。

\section{投資における時間の推移と収益率}

\subsection{価格変化と組み替えの考え方}

ポートフォリオ運用では，「価格変化」と「組み替え(リバランス)」が繰り返される。
2資産(債券価格 $B_t$，株価 $S_t$;保有量 $b_t$， $\Theta_t$)のケースで考えると，
時間の流れは
\[
\underbrace{t-1 \to t}_{\text{価格変化(保有量は不変)}}
\;\to\;
\underbrace{t}_{\text{組み替え}}
\;\to\;
\underbrace{t \to t+1}_{\text{価格変化(保有量は不変)}}
\;\to\;
\underbrace{t+1}_{\text{組み替え}}
\;\to\;\cdots
\]
となる。すなわち，時刻 $t$ の価格を参考にポートフォリオを組み替え，次の価格変化の
間は保有量は不変である。価格変化と組み替えは，同時刻で表示していても
“同時”ではないことに注意する(取引コストは無視する)。

\subsection{個別資産とポートフォリオの収益率}

\begin{itemize}
\item 資産数：$N$
\item 資産 $j$ の価格：$P_j(t)$(単位量当たりの価格)
\item 資産 $j$ の保有量：$x_j$(一期間の間で不変と考える)
\item ポートフォリオの価格：$V_P(t)$
\item 資産 $j$ の保有比率：$w_j(t) = x_j P_j(t)/V_P(t)$
\end{itemize}

このとき，ポートフォリオの(算術)収益率は
\begin{align}
R_P(t) &= \frac{V_P(t) - V_P(t-1)}{V_P(t-1)}
= \frac{1}{V_P(t-1)}\sum_{j=1}^{N} x_j\bigl(P_j(t) - P_j(t-1)\bigr)\notag\\
&= \sum_{j=1}^{N}\frac{x_j P_j(t-1)}{V_P(t-1)}\times
\frac{P_j(t)-P_j(t-1)}{P_j(t-1)}
= \sum_{j=1}^{N} w_j(t-1)\, R_j(t)
\label{eq:port-return}
\end{align}
となる。すなわち，\textbf{ポートフォリオの収益率は，構成資産の収益率の保有比率に
よる加重平均}である。ここで $R_j(t)$ と $R_P(t)$ は算術収益率である。対数収益率
では，このようにうまく表現できない(これが，ポートフォリオの最適化で算術収益率が
使われる理由である)。

\section{分散投資効果の例}

リスクを標準偏差(または分散)で測り，分散投資の効果を2つのケースで確認する。
いずれも，ある金額(例えば100億円)を持っている人が投資先を増やしていく
(1件当たりの投資額は減っていく)ことを想定する。

\subsection{ケース1：独立な収益率}

$N$ 銘柄の資産への投資について，次を仮定する。
\begin{enumerate}
\item すべての投資先の収益率 $R_j$， $j=1,\dots,N$ は，同じ分布に従う
($\E[R_j]=\mu$， $\Var[R_j]=\sigma^2$)。
\item すべての投資先の収益率 $R_j$ は互いに独立である。
\item すべての投資先への投資比率(保有比率)は同じとする。投資比率の和は1
なので，$w_j = 1/N$ となる。
\item ポートフォリオの収益率は，全投資先の収益率の(投資比率による)加重平均
\eqref{eq:port-return}で与えられる。
\item リスクは収益率の標準偏差で測る。
\end{enumerate}

このとき，ポートフォリオの収益率は
\[
R_p = \sum_{j=1}^{N} w_j R_j = \frac{1}{N}\sum_{j=1}^{N} R_j
\]
であり，期待値と分散は
\begin{align}
\E[R_p] &= \E\Bigl[\frac{1}{N}\sum_{j=1}^{N}R_j\Bigr]
= \frac{1}{N}\sum_{j=1}^{N}\E[R_j] = \frac{1}{N}\sum_{j=1}^{N}\mu = \mu,\\
\Var[R_p] &= \Var\Bigl[\frac{1}{N}\sum_{j=1}^{N}R_j\Bigr]
= \frac{1}{N^2}\Var\Bigl[\sum_{j=1}^{N}R_j\Bigr]
= \frac{1}{N^2}\sum_{j=1}^{N}\Var[R_j]
= \frac{\sigma^2}{N}
\end{align}
となる。ここで，期待値の計算の等号は常に成り立つ(期待値の線形性)が，分散の
計算で「和の分散 $=$ 分散の和」とした等号は常に成り立つものではなく，
\emph{無相関だから}成り立つことに注意する。したがって，ポートフォリオのリスクは
\begin{equation}
\sigma_p \equiv \sqrt{\Var[R_p]} = \frac{\sigma}{\sqrt{N}}
\;\longrightarrow\; 0 \qquad (N\to\infty).
\end{equation}
投資先の数 $N$ を増やすとリスクは $1/\sqrt{N}$ のオーダーで減少し，ゼロに
漸近する。

\subsection{ケース2：共通成分がある収益率}

ケース1の仮定1， 2を次のように変更する(他の仮定はそのまま)。

\begin{itemize}
\item[1'.] 投資先の収益率 $R_j$， $j=1,\dots,N$ は
\begin{equation}
R_j = \beta R + \varepsilon_j
\end{equation}
で与えられる。ただし，$R$ と $\varepsilon_j$ はそれぞれ互いに独立な確率変数で，
$\varepsilon_j$ と $\varepsilon_k$($j\neq k$)も独立とする。また
\[
\E[R] = \mu,\quad \Var[R] = \sigma^2,\qquad
\E[\varepsilon_j] = 0,\quad \Var[\varepsilon_j] = \sigma_\varepsilon^2.
\]
\end{itemize}

このとき，
\[
R_p = \sum_{j=1}^{N}\frac{1}{N}R_j
= \beta R + \frac{1}{N}\sum_{j=1}^{N}\varepsilon_j
\]
なので，
\begin{align}
\E[R_p] &= \beta\E[R] + \frac{1}{N}\sum_{j=1}^{N}\E[\varepsilon_j] = \beta\mu,\\
\Var[R_p] &= \beta^2\Var[R] + \frac{1}{N^2}\sum_{j=1}^{N}\Var[\varepsilon_j]
= \beta^2\sigma^2 + \frac{\sigma_\varepsilon^2}{N}
\end{align}
となる(ここでも分散の等号は独立性による)。よって
\begin{equation}
\sigma_p = \sqrt{\beta^2\sigma^2 + \frac{\sigma_\varepsilon^2}{N}}
\;\longrightarrow\; \beta\sigma \qquad (N\to\infty).
\end{equation}

\subsection{システマティック・リスクとアンシステマティック・リスク}

\begin{keypoint}[分散投資効果]
\begin{itemize}
\item リスクのうち，「投資先ごとに独立な成分」(\textbf{アンシステマティック・
リスク})は，投資先を分散させる($N$を増やす)につれて低下し，$N\to\infty$ では
ゼロに漸近する。
\item リスクのうち，「投資先に共通な成分」(\textbf{システマティック・リスク})は，
いくら投資先を分散させても低下しない。「投資先に共通な成分」の大きさを示すのが
$\beta$ である。
\end{itemize}
このように，投資先を分散することでリスクが低減することを\textbf{分散投資効果}と
いう。ここで示した例は分散投資効果の一端であり，実はより幅広い効果がある
(次節以降)。
\end{keypoint}

\section{個別資産の収益率とポートフォリオ}

本節では，少数の将来シナリオを使った説明を行う(前節とは切り離して考える)。

\subsection{ケース1：2つのシナリオ}

\begin{example}[基礎データ]
将来のシナリオI， II(発生確率各50\%)に対する3資産の収益率を次のように設定する。
\vspace{0.9em}
\begin{center}
\begin{tabular}{lcccc}
\toprule
 & I & II & 期待値 & 標準偏差\\
\midrule
発生確率 & 50\% & 50\% & &\\
資産A & 30\% & 10\% & 20\% & 10.0\%\\
資産B & 50\% & 10\% & 30\% & 20.0\%\\
資産C & 10\% & 50\% & 30\% & 20.0\%\\
\bottomrule
\end{tabular}
\end{center}
\vspace{0.9em}
資産Bと資産Cは，期待値と標準偏差(リスク)でみると区別できない(無差別である)。
このとき，「Aが50\%，Bが50\%」のポートフォリオ$\alpha$と，「Aが50\%，Cが50\%」
のポートフォリオ$\beta$の，リターンとリスクはどうなるか。
\end{example}

計算の内訳を$\alpha$で示す。収益率(確率変数)を $R_A, R_B$ とすると，
ポートフォリオ$\alpha$の収益率は(算術収益率で)
\[
R_\alpha = \frac{R_A + R_B}{2}
\qquad\Bigl(R_\alpha = \sum_i w_i R_i,\ \sum_i w_i = 1\Bigr)
\]
であり，リターンは
\[
\E[R_\alpha] = 0.5\times\frac{0.3+0.5}{2} + 0.5\times\frac{0.1+0.1}{2} = 0.25,
\]
収益率の分散($=$リスク$^2$)は
\[
\Var[R_\alpha] = 0.5\Bigl(\frac{0.3+0.5}{2}-0.25\Bigr)^2
+ 0.5\Bigl(\frac{0.1+0.1}{2}-0.25\Bigr)^2 = 0.15^2.
\]
ポートフォリオ$\beta$も同様に計算すると，結果は次のとおりである。

\vspace{0.9em}
\begin{center}
\begin{tabular}{lcccc}
\toprule
 & I & II & 期待値 & 標準偏差\\
\midrule
発生確率 & 50\% & 50\% & &\\
$\alpha$(A50\%+B50\%) & 40\% & 10\% & 25\% & 15.0\%\\
$\beta$(A50\%+C50\%) & 20\% & 30\% & 25\% & 5.0\%\\
\bottomrule
\end{tabular}
\end{center}
\vspace{0.9em}

\begin{itemize}
\item リターンは同じだが，リスク(収益率の標準偏差)は違う。
\item ポートフォリオ$\alpha$のリスクは，AとBのリスクの平均値と同じである。
\item リスク回避的な人ならば，ポートフォリオ$\beta$を選択する。
\end{itemize}

構成比率を変えても同じことが起きる。例えば A 20\%・B(C) 80\% では
$\alpha$：期待値28\%・標準偏差18\%，$\beta$：期待値28\%・標準偏差14\%となり，
A 80\%・B(C) 20\% では$\alpha$：期待値22\%・標準偏差12\%，$\beta$：期待値22\%・
標準偏差4\%となる。構成比率を連続的に変えて(リスク，期待収益率)平面上に軌跡を
描くと，$\alpha$(A+B)は直線を描く一方，$\beta$(A+C)は左側に折れ込む軌跡と
なり，リターンが同じ場合，ポートフォリオ$\beta$の方が常にリスクが低い。
特にA+Cでは，ある比率でリスクがちょうど0になる点も現れる。

\subsection{なぜリスク量が変わるのか}

\begin{itemize}
\item 資産Aと資産Bは，収益率が\textbf{同じ方向}に動く(Aの収益率が高いときはBも
高く，低いときはBも低い)。そのためポートフォリオ$\alpha$は，良い時は良いが，
悪い時は悪い(収益率の変動性が高い)。
\item 資産Aと資産Cは，収益率が\textbf{逆の方向}に動く。そのためポートフォリオ
$\beta$は，各資産がお互いにカバーしあうので，比較的安定した収益が得られる
(収益率の変動性が低い)。
\end{itemize}

この「同時的な動き方」を示す尺度が，\textbf{相関係数}である。

\subsection{共分散と相関係数}

\begin{definition}[共分散と相関係数]
確率変数 $X, Y$ に対して，\textbf{共分散}を
\begin{equation}
\Cov[X,Y] = C[X,Y] = \E\bigl[(X-\E[X])(Y-\E[Y])\bigr],
\end{equation}
\textbf{相関係数}を
\begin{equation}
\Corr[X,Y] = \rho[X,Y]
= \frac{\Cov[X,Y]}{\sqrt{\Var[X]}\sqrt{\Var[Y]}}
= \E\Biggl[\Bigl(\frac{X-\E[X]}{\sqrt{\Var[X]}}\Bigr)
\Bigl(\frac{Y-\E[Y]}{\sqrt{\Var[Y]}}\Bigr)\Biggr]
\end{equation}
と定義する。相関係数は標準偏差でスケーリングされた共分散であり，
\begin{equation}
-1 \leq \rho[X,Y] \leq 1
\end{equation}
を満たす。
\end{definition}

\vspace{18pt}

\subsection{ケース2：4つのシナリオ}

\begin{example}
シナリオを4つに増やし，資産Aとの相関係数が異なる3つの資産B， C， Dを考える。
\vspace{0.9em}
\begin{center}
\small
\begin{tabular}{lccccccc}
\toprule
 & I & II & III & IV & 平均 & 標準偏差 & 相関係数\footnotemark\\
\midrule
発生確率 & 25\% & 25\% & 25\% & 25\% & & &\\
資産A & 30\% & 30\% & 10\% & 10\% & 20.0\% & 10.00\% & 1.0\\
資産B & 50\% & 50\% & 10\% & 10\% & 30.0\% & 20.00\% & 1.0\\
資産C & 10\% & 10\% & 50\% & 50\% & 30.0\% & 20.00\% & $-1.0$\\
資産D & 50\% & 10\% & 50\% & 10\% & 30.0\% & 20.00\% & 0.0\\
\bottomrule
\end{tabular}
\end{center}
\vspace{0.9em}
\footnotetext{資産Aの収益率との相関係数。}
資産AとB， C， Dをそれぞれ組み合わせて構成比率を変えると，(リスク，期待収益率)
平面上に，A+Bは直線，A+Cはリスク0の点を通る折れ線，A+Dはその中間の
左に膨らんだ曲線の軌跡を描く。
\end{example}

\begin{keypoint}
リターンが同じ場合，\textbf{相関係数の低い資産と組み合わせるほど，ポートフォリオの
リスクは小さくなる}。構成比率を動かして得られる点の集合を\textbf{投資機会集合}と
呼ぶ。なお，相関係数が $\pm 1$ になることは稀(普通はない)ので，A+Dのような
カーブになるのが一般的である。
\end{keypoint}

\subsection{相関とリスクの関係：2資産での確認}

\textbf{設定。}資産Aの収益率 $\tilde R_A$，資産Bの収益率 $\tilde R_B$ について
\[
\E[\tilde R_A] = \mu_A,\ \Var[\tilde R_A] = \sigma_A^2,\qquad
\E[\tilde R_B] = \mu_B,\ \Var[\tilde R_B] = \sigma_B^2,\qquad
\rho[\tilde R_A, \tilde R_B] = \rho_{AB}
\]
とし，資産Aへの投資比率を $x$($0\leq x\leq 1$)，ポートフォリオの収益率を
$\tilde R_P = x\tilde R_A + (1-x)\tilde R_B$ とする。

\textbf{ポートフォリオの収益率の分散。}
\begin{align}
\sigma_P^2 = \Var[\tilde R_P]
&= \Var[x\tilde R_A + (1-x)\tilde R_B]\notag\\
&= x^2\sigma_A^2 + 2x(1-x)\sigma_A\sigma_B\rho_{AB} + (1-x)^2\sigma_B^2\notag\\
&\leq x^2\sigma_A^2 + 2x(1-x)\sigma_A\sigma_B + (1-x)^2\sigma_B^2
= \bigl(x\sigma_A + (1-x)\sigma_B\bigr)^2,
\end{align}
ここで不等号は $\rho_{AB}\leq 1$ による。

\textbf{ポートフォリオの収益率の標準偏差。}
\begin{equation}
\sigma_P = \sqrt{x^2\sigma_A^2 + 2x(1-x)\sigma_A\sigma_B\rho_{AB}
+ (1-x)^2\sigma_B^2}
\;\leq\; x\sigma_A + (1-x)\sigma_B.
\end{equation}
等号が成り立つのは $\rho_{AB}=1$ のときで，そのとき
$\sigma_P = x\sigma_A + (1-x)\sigma_B$，すなわちポートフォリオのリスクは各資産の
リスクの加重平均値になる。

\subsection{一般の $n$ 資産への拡張}

\begin{keypoint}[相関係数とポートフォリオのリスク]
$x_j$ を資産 $j$ への投資比率($\sum_{j=1}^n x_j = 1$)とする。
\begin{itemize}
\item 全資産の収益率が「相関係数 $=1$」の場合：
\[
\sigma_P = \sum_{j=1}^{n} x_j\sigma_j \quad(\text{構成資産のリスクの加重平均}).
\]
\item それ以外の(1つでも相関係数 $\neq 1$ の資産がある)場合：
\[
\sigma_P < \sum_{j=1}^{n} x_j\sigma_j.
\]
\end{itemize}
一方，収益率の期待値は必ず加重平均値になる：
$\mu_P = \sum_{j=1}^{n} x_j\mu_j$(期待値の線形性より)。
\end{keypoint}

すなわち，\textbf{分散投資効果}により，さまざまな資産に投資することで，
ポートフォリオのリスクは(各構成資産のリスクの加重平均値よりも)小さくなる。
ポートフォリオ選択問題では，このような分散投資効果(分散投資によるリスクの
低減効果)をトコトン追求して，自身にとって最適なポートフォリオを導出する。

\subsection{どのポートフォリオを選ぶべきか：無差別曲線}

投資機会集合の中から，自分にとって最も望ましいポートフォリオ，すなわち
\textbf{自分の期待効用を最大にするポートフォリオ}を選ぶ。実際には数値的に
求めればよいが，ここではグラフィカルに考える。

\paragraph{無差別曲線の概形：指数効用の場合。}
\textbf{無差別曲線}とは，(ここでは)期待効用が一定値となるカーブのことである。
\begin{enumerate}
\item 指数効用関数 $u(w) = -e^{-aw}$($a>0$)を使用する。
\item $W \sim N(\mu, \sigma^2)$ を仮定する(ここでは平均と標準偏差で評価する
ので，正規分布を想定する)。
\end{enumerate}
すると，正規分布のモーメント母関数より，期待効用は
\begin{equation}
\E[u(W)] = -\E[e^{-aW}] = -e^{-a\E[W] + \frac{a^2}{2}\Var[W]}
= -e^{-a\mu + \frac{a^2}{2}\sigma^2}
\end{equation}
となる。期待効用 $=$ 一定となるとき，$-a\mu + \frac{a^2}{2}\sigma^2 = $ 一定，
すなわち
\[
\mu = \frac{a}{2}\sigma^2 + \text{定数}
\]
となり，\textbf{$\mu$ は $\sigma$ の二次関数(かつ $\sigma$ の二次項の係数は正)}
である。つまり$(\sigma,\mu)$平面上で，無差別曲線は下に凸の右上がりの放物線で
あり，左上にある無差別曲線ほど期待効用は高い。
(これだけでは一般性はないが，例として。)

\paragraph{最適ポートフォリオの選択。}
$(\sigma,\mu)$平面上に，期待効用の値 $u_1 < u_2 < u_3$ に対応する無差別曲線群と
投資機会集合を重ねて描くと，期待効用が最大になるのは，
\textbf{ポートフォリオの投資機会集合と無差別曲線の接点}である。この接点の
ポートフォリオが(この投資家にとっての)最適ポートフォリオとなる。

\begin{remark}[無リスクポートフォリオの存在について]
相関係数 $=-1$ の資産の組み合わせで無リスクポートフォリオができることを示す図が
あった。数理的には，図に示したように無リスクポートフォリオを作れる。そして，
無リスクポートフォリオを作ったら，その価格は無リスク資産と同一である(後で紹介
するとおり，無リスク資産は市場に唯一つであり，価格の異なる複数の無リスク資産が
市場に存在してはいけない。これはデリバティブの価格付けの考え方の基礎にもなって
いる)。そのため，この無リスクポートフォリオは無リスク資産そのものになる。

なお，普通は，相関係数 $=-1$ となる\emph{異なる}資産は存在しない。ある資産と
完全に真逆の動きをする資産は，もとの資産の空売りである。そして，資産の保有量と
同量を空売りすると，結局何も投資していないことになる。
\end{remark}

\vspace{18pt}

\section{補足：最小分散ポートフォリオ(Markowitz)}

本節は，配布資料「最小分散ポートフォリオ(Markowitz)」に基づき，一期間の
平均分散モデルによる最小分散ポートフォリオの算出などを扱う。ここではリスク性
資産のみからなる場合(Markowitzのポートフォリオ選択問題)を扱う。

\subsection{設定}

$N$ 個のリスク性資産からなるポートフォリオを考える。資産 $i$($i=1,\dots,N$)の
収益率を確率変数 $R_i$，$\bm{R} = (R_1,\dots,R_N)^\top$ を $N$ 次元確率変数
ベクトルとし，$\bm{R}$ の期待値ベクトル
$\bm{\mu} = (\mu_1,\dots,\mu_N)^\top = (\E[R_1],\dots,\E[R_N])^\top$ と
分散共分散行列 $\bm{\Sigma} = (\sigma_{ij})_{i,j=1,\dots,N}$，
$\sigma_{ij} = \Cov(R_i,R_j)$ は既知とする。$\bm{e} = (1,\dots,1)^\top$ とすると，
平均分散モデルにおける\textbf{最小分散ポートフォリオ}
$\bm{x} = (x_1,\dots,x_N)^\top$($x_i$ は資産 $i$ への投資比率)は
\begin{equation}
\hat{\bm{x}} = \argmin_{\bm{x}\in\mathcal{R}^N}
\frac{1}{2}\bm{x}^\top\bm{\Sigma}\bm{x}
\label{eq:mk-problem}
\end{equation}
で与えられる。ただし，$\alpha$ を定数として，制約条件
\begin{align}
\bm{\mu}^\top\bm{x} &= \alpha,
\label{eq:mk-c1}\\
\bm{e}^\top\bm{x} &= 1
\label{eq:mk-c2}
\end{align}
を満たし，$N$ 次実対称行列 $\bm{\Sigma}$ は正定値と仮定する。また，$\bm{\mu}$ と
$\bm{e}$ は線形独立と仮定する。

\begin{remark}
\eqref{eq:mk-c1}はポートフォリオの収益率の期待値が $\alpha$ であること，
\eqref{eq:mk-c2}は $\bm{x}$ が投資比率であることを示している。
\end{remark}

\vspace{18pt}

\subsection{最小分散ポートフォリオの導出}

Lagrange未定乗数法を用いる。Lagrange未定乗数を $\lambda_1, \lambda_2$ として，
Lagrange関数を
\begin{equation}
L = \frac{1}{2}\bm{x}^\top\bm{\Sigma}\bm{x}
- \lambda_1(\bm{\mu}^\top\bm{x} - \alpha)
- \lambda_2(\bm{e}^\top\bm{x} - 1)
\end{equation}
とすると，求める解 $\hat{\bm{x}}$ は
\begin{align}
\frac{\partial L}{\partial x_i}
&= \frac{\partial}{\partial x_i}\Biggl[\frac{1}{2}\sum_{j,k=1}^{N}
\sigma_{jk}x_j x_k - \lambda_1\Bigl(\sum_{j=1}^N \mu_j x_j - \alpha\Bigr)
- \lambda_2\Bigl(\sum_{j=1}^N x_j - 1\Bigr)\Biggr]\notag\\
&= \sum_{k=1}^{N}\sigma_{ik}x_k - \lambda_1\mu_i - \lambda_2 = 0,
\qquad i=1,\dots,N,
\label{eq:mk-foc}\\
\frac{\partial L}{\partial\lambda_1}
&= -(\bm{\mu}^\top\bm{x} - \alpha) = 0,\\
\frac{\partial L}{\partial\lambda_2}
&= -(\bm{e}^\top\bm{x} - 1) = 0
\end{align}
を満たす。\eqref{eq:mk-foc}をベクトル表記で
\begin{equation}
\frac{\partial L}{\partial\bm{x}}
= \bm{\Sigma}\bm{x} - \lambda_1\bm{\mu} - \lambda_2\bm{e} = 0
\end{equation}
と書き直すと，$\bm{\Sigma}$ は正定値と仮定したので逆行列 $\bm{\Sigma}^{-1}$ が
存在し，最小分散ポートフォリオは
\begin{equation}
\hat{\bm{x}} = \bm{\Sigma}^{-1}(\lambda_1\bm{\mu} + \lambda_2\bm{e})
\label{eq:mk-x-lambda}
\end{equation}
で与えられる。ただし，$\lambda_1,\lambda_2$ が未知変数なので，
\eqref{eq:mk-x-lambda}はまだ具体的な最小分散ポートフォリオを示していない。

そこで，未知数 $\lambda_1,\lambda_2$ を求める。\eqref{eq:mk-x-lambda}を制約条件に
代入すると，
\begin{align}
\bm{\mu}^\top\hat{\bm{x}} - \alpha
&= \lambda_1\bm{\mu}^\top\bm{\Sigma}^{-1}\bm{\mu}
+ \lambda_2\bm{\mu}^\top\bm{\Sigma}^{-1}\bm{e} - \alpha = 0,\\
\bm{e}^\top\hat{\bm{x}} - 1
&= \lambda_1\bm{e}^\top\bm{\Sigma}^{-1}\bm{\mu}
+ \lambda_2\bm{e}^\top\bm{\Sigma}^{-1}\bm{e} - 1 = 0
\end{align}
となる。ここで
\begin{equation}
a = \bm{\mu}^\top\bm{\Sigma}^{-1}\bm{\mu},\qquad
b = \bm{\mu}^\top\bm{\Sigma}^{-1}\bm{e} = \bm{e}^\top\bm{\Sigma}^{-1}\bm{\mu},\qquad
c = \bm{e}^\top\bm{\Sigma}^{-1}\bm{e}
\label{eq:mk-abc}
\end{equation}
とおくと，上の2式は
\begin{equation}
\begin{pmatrix} a & b\\ b & c\end{pmatrix}
\begin{pmatrix}\lambda_1\\ \lambda_2\end{pmatrix}
= \begin{pmatrix}\alpha\\ 1\end{pmatrix}
\end{equation}
と書けるので，$d = ac - b^2$ として
\begin{equation}
\begin{pmatrix}\lambda_1\\ \lambda_2\end{pmatrix}
= \frac{1}{d}\begin{pmatrix} c & -b\\ -b & a\end{pmatrix}
\begin{pmatrix}\alpha\\ 1\end{pmatrix}
\end{equation}
が得られる\footnote{$\bm{\mu}$ と $\bm{e}$ は線形独立なので
$d = ac - b^2 \neq 0$ が成り立つ。もしも線形従属ならば，任意の $\alpha$ に対して
$(a,b) = \alpha(b,c)$ でなければならない。もしも $\bm{\Sigma}$ が正定値であれば
$a>0$ かつ $c>0$，さらに
$(b\bm{\mu}-a\bm{e})^\top\bm{\Sigma}^{-1}(b\bm{\mu}-a\bm{e}) = ad > 0$ より
$d>0$ となる。ここでは $\bm{\Sigma}$ を正定値と仮定したが，$\bm{\Sigma}$ は必ず
実対称行列で非負定値(半正定値)になるものの，正定値になるとは限らない。
正定値行列は正則，すなわち逆行列をもつ。}。これを\eqref{eq:mk-x-lambda}に
代入すると，具体的な最小分散ポートフォリオ
\begin{equation}
\hat{\bm{x}} = \hat{\bm{x}}(\alpha)
= \bm{\Sigma}^{-1}\Bigl(\frac{c\alpha - b}{d}\bm{\mu}
+ \frac{a - b\alpha}{d}\bm{e}\Bigr)
= \bm{g} + \alpha\bm{h}
= (1-\alpha)\bm{g} + \alpha(\bm{g}+\bm{h})
\label{eq:mk-solution}
\end{equation}
が得られる。ただし，
\begin{equation}
\bm{g} = \frac{a\bm{\Sigma}^{-1}\bm{e} - b\bm{\Sigma}^{-1}\bm{\mu}}{d},
\qquad
\bm{h} = \frac{c\bm{\Sigma}^{-1}\bm{\mu} - b\bm{\Sigma}^{-1}\bm{e}}{d}
\label{eq:mk-gh}
\end{equation}
である。

\subsection{最小分散ポートフォリオの解の意味}

\eqref{eq:mk-solution}より，最小分散ポートフォリオ $\hat{\bm{x}}$ は $\bm{g}$ と
$\bm{h}$ の線形結合(あるいは $\bm{g}$ と $\bm{g}+\bm{h}$ の分点)として表現
される。\eqref{eq:mk-abc}より，$\bm{g}$ と $\bm{h}$ の分子はそれぞれ
\begin{align}
\bm{\mu}^\top(a\bm{\Sigma}^{-1}\bm{e} - b\bm{\Sigma}^{-1}\bm{\mu})
&= ab - ba = 0,\\
\bm{\mu}^\top(c\bm{\Sigma}^{-1}\bm{\mu} - b\bm{\Sigma}^{-1}\bm{e})
&= ca - b^2 = d
\end{align}
を満たすので，
\begin{equation}
\bm{\mu}^\top\bm{g} = 0,\qquad \bm{\mu}^\top\bm{h} = 1
\end{equation}
が得られる。これより，
\begin{itemize}
\item $\bm{g}$ は期待収益率0のポートフォリオ，
\item $\bm{h}$(および $\bm{g}+\bm{h}$)は期待収益率1のポートフォリオ，
\end{itemize}
である。最小分散ポートフォリオは $\bm{g}$ と $\bm{h}$ の線形結合
\eqref{eq:mk-solution}で表現され，二つのポートフォリオ $\bm{g}$ と
$\bm{h}'\equiv\bm{g}+\bm{h}$ の混合比率はポートフォリオの期待収益率 $\alpha$ に
依存する。また，期待収益率 $\mu$ を与えると最小分散ポートフォリオ $\hat{\bm{x}}$
は一意に決まることもわかる。

\subsection{最小分散境界}

次に，\textbf{最小分散境界}を求める。\eqref{eq:mk-solution}，\eqref{eq:mk-gh}
より，最小分散ポートフォリオの分散 $\hat\sigma^2$ は
\begin{align}
\hat\sigma^2 = \Var[\bm{R}^\top\hat{\bm{x}}]
= \hat{\bm{x}}^\top\bm{\Sigma}\hat{\bm{x}}
&= (\bm{g}+\alpha\bm{h})^\top\bm{\Sigma}(\bm{g}+\alpha\bm{h})\notag\\
&= \bm{g}^\top\bm{\Sigma}\bm{g}
+ (\bm{h}^\top\bm{\Sigma}\bm{g} + \bm{g}^\top\bm{\Sigma}\bm{h})\alpha
+ \bm{h}^\top\bm{\Sigma}\bm{h}\,\alpha^2
\end{align}
である。\eqref{eq:mk-abc}，\eqref{eq:mk-gh}などより，
\begin{equation}
\bm{g}^\top\bm{\Sigma}\bm{g} = \frac{a}{d},\qquad
\bm{h}^\top\bm{\Sigma}\bm{g} = \bm{g}^\top\bm{\Sigma}\bm{h} = -\frac{b}{d},\qquad
\bm{h}^\top\bm{\Sigma}\bm{h} = \frac{c}{d}
\end{equation}
となるので，これらを代入すると
\begin{equation}
\hat\sigma^2 = \frac{a - 2b\alpha + c\alpha^2}{d}
\label{eq:mk-frontier}
\end{equation}
が得られる。あるいは，\eqref{eq:mk-frontier}を二次曲線の標準形に変形すると，
\begin{equation}
\frac{\hat\sigma^2}{\bigl(\frac{1}{\sqrt{c}}\bigr)^2}
- \frac{\bigl(\alpha - \frac{b}{c}\bigr)^2}{\bigl(\frac{\sqrt{d}}{c}\bigr)^2} = 1,
\qquad \hat\sigma > 0
\label{eq:mk-hyperbola}
\end{equation}
が得られる。\eqref{eq:mk-hyperbola}は，最小分散境界(すなわち最小分散
ポートフォリオのリスク $\hat\sigma$ と期待収益率 $\alpha$)が
$(\sigma,\alpha)$ 平面上で\textbf{双曲線}を描くことを示している。これは，
二次曲線の分類(本節末尾の付録A)における判別式が
\begin{equation}
\delta \equiv \det\bm{Q} = -c\times\frac{c^2}{ac-b^2} = -\frac{c^3}{d} < 0
\end{equation}
となることからわかる。

\subsection{大域的最小分散ポートフォリオ}

\eqref{eq:mk-hyperbola}より，分散が最小になるのは期待収益率が
$\alpha^* = b/c$ のときで，このときの分散は $(\sigma^*)^2 = 1/c$，そして
\eqref{eq:mk-solution}より\textbf{大域的最小分散ポートフォリオ}
(global minimum variance portfolio) $\hat{\bm{x}}^*$ は
\begin{equation}
\hat{\bm{x}}^* = \bm{g} + \frac{b}{c}\bm{h}
= \frac{ac\,\bm{\Sigma}^{-1}\bm{e} - b^2\bm{\Sigma}^{-1}\bm{e}}{dc}
= \frac{1}{c}\bm{\Sigma}^{-1}\bm{e}
\end{equation}
で与えられる。

\subsection{最小分散ポートフォリオの性質}

\subsubsection{二資産分離定理}

Markowitzの問題でも次の\textbf{二資産分離定理}が成り立つ。

\begin{theorem}[二資産分離定理]
最小分散ポートフォリオは，任意の2つの異なる最小分散ポートフォリオの線形結合で
表現できる。
\label{thm:two-fund}
\end{theorem}

このことは，$\alpha_1\neq\alpha_2$ として，期待収益率 $\mu_1$ の最小分散
ポートフォリオ $\hat{\bm{x}}_1$ と期待収益率 $\mu_2$ の最小分散ポートフォリオ
$\hat{\bm{x}}_2$ を選ぶと，
\[
\hat{\bm{x}}_1 = \bm{g} + \mu_1\bm{h},\qquad
\hat{\bm{x}}_2 = \bm{g} + \mu_2\bm{h}
\]
と書けるので，期待収益率 $\mu$ の最小分散ポートフォリオ $\hat{\bm{x}}(\mu)$ が
\begin{equation}
\hat{\bm{x}}(\mu) = \frac{\mu_2 - \mu}{\mu_2 - \mu_1}\hat{\bm{x}}_1
+ \frac{\mu - \mu_1}{\mu_2 - \mu_1}\hat{\bm{x}}_2
\end{equation}
と表現できることから明らかである。

\subsubsection{ポートフォリオ間の共分散}

\begin{theorem}
期待収益率 $\mu_1$ の最小分散ポートフォリオ $\hat{\bm{x}}_1$ と期待収益率
$\mu_2$ の最小分散ポートフォリオ $\hat{\bm{x}}_2$ の収益率をそれぞれ
$R(\hat{\bm{x}}_1), R(\hat{\bm{x}}_2)$ とすると，その共分散は
\begin{equation}
C\bigl(R(\hat{\bm{x}}_1), R(\hat{\bm{x}}_2)\bigr)
= \frac{c}{d}\Bigl(\mu_1 - \frac{b}{c}\Bigr)\Bigl(\mu_2 - \frac{b}{c}\Bigr)
+ \frac{1}{c}
\label{eq:mk-cov}
\end{equation}
で与えられる。
\label{thm:mk-cov}
\end{theorem}

\begin{proof}
$\hat{\bm{x}}_1 = \bm{g}+\mu_1\bm{h}$， $\hat{\bm{x}}_2 = \bm{g}+\mu_2\bm{h}$ と
書けるので，
\begin{align*}
C\bigl(R(\hat{\bm{x}}_1), R(\hat{\bm{x}}_2)\bigr)
&= \hat{\bm{x}}_1^\top\bm{\Sigma}\hat{\bm{x}}_2
= (\bm{g}+\mu_1\bm{h})^\top\bm{\Sigma}(\bm{g}+\mu_2\bm{h})\\
&= \frac{a}{d} + (\mu_1+\mu_2)\Bigl(-\frac{b}{d}\Bigr) + \mu_1\mu_2\frac{c}{d}
= \frac{c}{d}\mu_1\mu_2 - (\mu_1+\mu_2)\frac{c}{d}\cdot\frac{b}{c}
+ \frac{b^2}{dc} - \frac{b^2}{dc} + \frac{a}{d}\\
&= \frac{c}{d}\Bigl(\mu_1-\frac{b}{c}\Bigr)\Bigl(\mu_2-\frac{b}{c}\Bigr)
+ \frac{ac-b^2}{dc}
= \frac{c}{d}\Bigl(\mu_1-\frac{b}{c}\Bigr)\Bigl(\mu_2-\frac{b}{c}\Bigr)
+ \frac{1}{c}.\qedhere
\end{align*}
\end{proof}

特に，$\hat{\bm{x}}_1 = \hat{\bm{x}}_2 = \hat{\bm{x}}$ とすると，
\begin{equation}
V[R(\hat{\bm{x}})]
= \frac{c}{d}\Bigl(\E[R(\hat{\bm{x}})] - \frac{b}{c}\Bigr)^2 + \frac{1}{c}
\label{eq:mk-var}
\end{equation}
より，次の定理が得られる。

\begin{theorem}
任意の最小分散ポートフォリオ $\hat{\bm{x}}$ において，
\begin{equation}
V[R(\hat{\bm{x}})] \geq \frac{1}{c}
\end{equation}
が成り立つ。等号は $\hat{\bm{x}} = \hat{\bm{x}}^*$(大域的最小分散
ポートフォリオ)でのみ成り立つ。
\end{theorem}

\vspace{18pt}

\subsubsection{ゼロベータポートフォリオ}

定理\ref{thm:mk-cov}で，$\hat{\bm{x}}_1(\mu_1) = \hat{\bm{x}}(\mu)$，
$\hat{\bm{x}}(\mu) \neq \hat{\bm{x}}^*$ を固定して関数
$f(\mu_2) = C\bigl(R(\hat{\bm{x}}), R(\hat{\bm{x}}(\mu_2))\bigr)$ を考えると，
$f(\mu_2)$ は単調関数なので，$f(\mu_2) = 0$ を満たす $\mu_2 = \mu^o$ および
その最小分散ポートフォリオ $\hat{\bm{x}}^o \equiv \hat{\bm{x}}(\mu^o)$ は一意に
決まる。さらに $f(\mu^o) = C(R(\hat{\bm{x}}), R(\hat{\bm{x}}^o)) = 0$ より，
\eqref{eq:mk-cov}は
\begin{equation}
0 = \frac{c}{d}\Bigl(\mu-\frac{b}{c}\Bigr)\Bigl(\mu^o-\frac{b}{c}\Bigr)
+ \frac{1}{c}
\end{equation}
となるので，$\mu^o$ は $\mu$ を用いて
\begin{equation}
\mu^o = \frac{b}{c} - \frac{d}{c^2}\Bigl(\mu - \frac{b}{c}\Bigr)^{-1}
= \frac{b}{c} - \frac{d}{c(c\mu - b)}
\label{eq:mk-muo}
\end{equation}
と書ける。\eqref{eq:mk-muo}より，
\begin{equation}
\mu - \mu^o = \mu - \frac{b}{c} + \frac{d}{c(c\mu-b)}
= \frac{c(\mu-\frac{b}{c})^2 + \frac{d}{c}}{c\mu - b}
= d\,\frac{\frac{c}{d}(\mu-\frac{b}{c})^2 + \frac{1}{c}}{c\mu - b}
\end{equation}
となるが，\eqref{eq:mk-var}より
$V[R(\hat{\bm{x}})] = \frac{c}{d}(\mu-\frac{b}{c})^2 + \frac{1}{c}$ なので，
\begin{equation}
\mu - \mu^o = \frac{d}{c\mu - b}\,V[R(\hat{\bm{x}})]
\label{eq:mk-mu-muo}
\end{equation}
が得られる。また，定理\ref{thm:two-fund}より，任意の最小分散ポートフォリオ
$\hat{\bm{x}}'$ は $\hat{\bm{x}}$ と $\hat{\bm{x}}^o$ の線形結合で表現できるので，
実数 $\beta$ を用いて
\begin{equation}
\hat{\bm{x}}' = \beta\hat{\bm{x}} + (1-\beta)\hat{\bm{x}}^o
\end{equation}
と表現すると，
\begin{equation}
\E[R(\hat{\bm{x}}')]
= \bm{\mu}^\top\bigl(\beta(\bm{g}+\mu\bm{h}) + (1-\beta)(\bm{g}+\mu^o\bm{h})\bigr)
= \beta\mu + (1-\beta)\mu^o = \mu^o + \beta(\mu - \mu^o)
\end{equation}
と表現できる。一方，\eqref{eq:mk-cov}と\eqref{eq:mk-muo}より，
$\mu' \equiv \E[R(\hat{\bm{x}}')]$ は
\begin{equation}
\mu' = \frac{b}{c} - \frac{d}{c(c\mu-b)} + \frac{d}{c\mu-b}
C\bigl(R(\hat{\bm{x}}'),R(\hat{\bm{x}})\bigr)
= \mu^o + \frac{d}{c\mu-b}C\bigl(R(\hat{\bm{x}}'),R(\hat{\bm{x}})\bigr)
\end{equation}
となるので，\eqref{eq:mk-mu-muo}を用いて
\begin{equation}
\mu' = \mu^o + \frac{d}{c\mu-b}V[R(\hat{\bm{x}})]\times
\frac{C\bigl(R(\hat{\bm{x}}'),R(\hat{\bm{x}})\bigr)}{V[R(\hat{\bm{x}})]}
= \mu^o + (\mu-\mu^o)\times
\frac{C\bigl(R(\hat{\bm{x}}'),R(\hat{\bm{x}})\bigr)}{V[R(\hat{\bm{x}})]}.
\end{equation}
よって，次の定理が成り立つ。

\begin{theorem}
大域的最小分散ポートフォリオ $\hat{\bm{x}}^*$ を除く最小分散ポートフォリオ
$\hat{\bm{x}}$ に対して，
\begin{equation}
C\bigl(R(\hat{\bm{x}}), R(\hat{\bm{x}}^o)\bigr) = 0
\end{equation}
を満たす最小分散ポートフォリオ $\hat{\bm{x}}^o$ が存在し，一意に定まる。さらに，
任意の最小分散ポートフォリオ $\hat{\bm{x}}'$ は，$\hat{\bm{x}}$ と
$\hat{\bm{x}}^o$ の線形結合で表現できて，期待収益率は
\begin{equation}
\E[R(\hat{\bm{x}}')] = \beta\,\E[R(\hat{\bm{x}})] + (1-\beta)\E[R(\hat{\bm{x}}^o)],
\qquad
\beta = \frac{C\bigl(R(\hat{\bm{x}}'),R(\hat{\bm{x}})\bigr)}{V[R(\hat{\bm{x}})]}
\end{equation}
で与えられる。
\label{thm:zero-beta}
\end{theorem}

定理\ref{thm:zero-beta}の $\hat{\bm{x}}$ と $\hat{\bm{x}}^o$ は，一方が効率的
フロンティア上にあり，もう一方は効率的フロンティア上にない。これは，
\eqref{eq:mk-muo}より，$\mu$ と $\mu^o$ がともに $b/c$ を上回る($=$効率的
フロンティア上にある)ことはないことからわかる。$\hat{\bm{x}}^o$ を，
$\hat{\bm{x}}$ に関する\textbf{ゼロベータポートフォリオ} (zero-beta portfolio)
という。次の定理も重要である。

\begin{theorem}
大域的最小分散ポートフォリオ $\hat{\bm{x}}^*$ を除く最小分散ポートフォリオ
$\hat{\bm{x}}$ に関するゼロベータポートフォリオ $\hat{\bm{x}}^o$ の期待収益率は，
$(\sigma,\mu)$ 平面上の最小分散境界上の $\hat{\bm{x}}$ における接線と $\mu$ 軸との
切片に等しい。
\label{thm:zero-beta-tangent}
\end{theorem}

\begin{proof}
最小分散境界を表す双曲線\eqref{eq:mk-hyperbola}に陰関数定理を使うと，接線の
傾きは
\begin{equation}
m(\hat{\bm{x}}) = \frac{\dd \E[R(\hat{\bm{x}})]}{\dd V[R(\hat{\bm{x}})]}
\cdot\frac{\dd V[R(\hat{\bm{x}})]}{\dd\sigma[R(\hat{\bm{x}})]}
= \frac{d\,\sigma[R(\hat{\bm{x}})]}{c\,\E[R(\hat{\bm{x}})] - b}
\end{equation}
となるので，接線と $\mu$ 軸との切片は，\eqref{eq:mk-mu-muo}などより
\begin{align}
\E[R(\hat{\bm{x}})] - m(\hat{\bm{x}})\,\sigma[R(\hat{\bm{x}})]
&= \E[R(\hat{\bm{x}})] - \frac{d\,\sigma[R(\hat{\bm{x}})]}
{c\E[R(\hat{\bm{x}})]-b}\times\sigma[R(\hat{\bm{x}})]
= \E[R(\hat{\bm{x}})] - \frac{d\,V[R(\hat{\bm{x}})]}{c\E[R(\hat{\bm{x}})]-b}\notag\\
&= \E[R(\hat{\bm{x}})] - \frac{c(\mu-\frac{b}{c})^2 + \frac{d}{c}}
{c\E[R(\hat{\bm{x}})]-b}
= \E[R(\hat{\bm{x}}^o)]
\end{align}
となるが，これはゼロベータポートフォリオ $\hat{\bm{x}}^o$ の期待収益率に等しい
ことを示している。
\end{proof}

なお，ゼロベータポートフォリオは無リスク資産がない場合に拡張したCAPMの議論
(ゼロベータCAPMという)において使われる(本節末尾の付録B参照)。

\subsection*{付録A：二次曲線の分類(一部のみ)}

$a,b,c,d,e,h$ を定数とする二次形式
\begin{equation}
f(x,y) \equiv ax^2 + 2hxy + by^2 + 2dx + 2ey + c
\end{equation}
において，
\begin{equation}
\bm{A} \equiv \begin{pmatrix} a & h & d\\ h & b & e\\ d & e & c\end{pmatrix},
\qquad
\bm{Q} \equiv \begin{pmatrix} a & h\\ h & b\end{pmatrix}
\end{equation}
とおく。$\det\bm{A} \neq 0$ のとき，二次曲線 $f(x,y)=0$ の表す図形は，
$\det\bm{Q}$ の符号により
\begin{equation}
\begin{cases}
\det\bm{Q} > 0 & \longrightarrow\ \text{楕円}\\
\det\bm{Q} = 0 & \longrightarrow\ \text{放物線}\\
\det\bm{Q} < 0 & \longrightarrow\ \text{双曲線}
\end{cases}
\end{equation}
に分類される。典型的な方程式として，楕円は
$\frac{x^2}{a^2}+\frac{y^2}{b^2}=1$($a>0$， $b>0$)，放物線は $y=x^2$，
双曲線は $y = 1/x$ があり，それぞれの $\delta \equiv \det\bm{Q}$ の符号を確認
されたい。

\subsection*{付録B：ゼロベータCAPM}

定理\ref{thm:zero-beta-tangent}の最小分散ポートフォリオ $\hat{\bm{x}}^*$
(訂正：任意の最小分散ポートフォリオ $\hat{\bm{x}}$)に対応する $(\sigma,\mu)$
平面上の点を点M とし，任意のリスク性資産Aを示す点 $(\sigma_A,\mu_A)$ を点Aと
する。ただし，点Mと点Aは一致しないものとする。ポートフォリオ $\hat{\bm{x}}$ と
リスク性資産Aからなるポートフォリオ P が $(\sigma,\mu)$ 平面上に描くカーブを
双曲線AMと呼ぶと，双曲線AMは，定理\ref{thm:zero-beta-tangent}の接線および
リスク性資産からなるポートフォリオの最小分散境界と点Mで接する\footnote{双曲線AM
は必ず点Mを通るが，もしも点Mでリスク性資産からなるポートフォリオの最小分散境界と
交差するならば，最小分散ポートフォリオよりも効率的なポートフォリオが存在する
ことになるため，矛盾が生じる。}。これをもとに，以下では
\begin{enumerate}
\item 双曲線AMを求める。
\item 双曲線AMの接線の傾きを求める。
\item 点Mにおける曲線AMの接線が定理\ref{thm:zero-beta-tangent}の接線と一致する
ための条件を求める。
\end{enumerate}
の手順で議論する。

ポートフォリオPの収益率を $R_P$，$\mu_P = \E[R_P]$，
$\sigma_P = \sqrt{V[R_P]}$，リスク性資産Aの収益率を $R_A$，ポートフォリオPの
リスク性資産Aへの投資比率を $w$ とすると，
\begin{align}
R_P &= (1-w)R_M + wR_A,\\
\mu_P &= \E[R_P] = (1-w)\mu_M + w\mu_A,\\
\sigma_P^2 &= V[R_P] = (1-w)^2\sigma_M^2 + 2(1-w)w\sigma_{AM} + w^2\sigma_A^2,
\qquad \sigma_{AM} = C(R_M,R_A) = \rho_{AM}\sigma_M\sigma_A
\end{align}
であり，
\begin{equation}
\sigma_P = \sqrt{(1-w)^2\sigma_M^2 + 2(1-w)w\sigma_{AM} + w^2\sigma_A^2}
\end{equation}
である。

次に，双曲線AM上の点Pで接線を考えると，その傾きは
\begin{equation}
\frac{\dd\mu_P}{\dd\sigma_P} = \frac{\dd\mu_P/\dd w}{\dd\sigma_P/\dd w}
\end{equation}
であり，
\begin{equation}
\frac{\dd\mu_P}{\dd w} = \mu_A - \mu_M,
\end{equation}
また，
\begin{equation}
\frac{\dd\sigma_P^2}{\dd w} = 2\sigma_P\frac{\dd\sigma_P}{\dd w}
= -2(1-w)\sigma_M^2 + 2(1-2w)\sigma_{AM} + 2w\sigma_A^2
\end{equation}
より
\begin{equation}
\frac{\dd\sigma_P}{\dd w}
= \frac{-(1-w)\sigma_M^2 + (1-2w)\sigma_{AM} + w\sigma_A^2}{\sigma_P}
\end{equation}
なので，
\begin{equation}
\frac{\dd\mu_P}{\dd\sigma_P}
= \frac{\sigma_P(\mu_A - \mu_M)}
{-(1-w)\sigma_M^2 + (1-2w)\sigma_{AM} + w\sigma_A^2}
\end{equation}
が得られる。

ここで，点Mでは $w=0$ なので，点Mにおける双曲線AMの接線の傾きは
\begin{equation}
\left.\frac{\dd\mu_P}{\dd\sigma_P}\right|_{w=0}
= \frac{\sigma_M(\mu_A - \mu_M)}{\sigma_{AM} - \sigma_M^2}
\end{equation}
であるが，点Mにおける接線は1つに限られるので，この傾きは
定理\ref{thm:zero-beta-tangent}の接線の傾きと一致しなければならない。ゆえに，
\begin{equation}
\frac{\sigma_M(\mu_A-\mu_M)}{\sigma_{AM}-\sigma_M^2}
= \frac{\E[R(\hat{\bm{x}})] - \E[R(\hat{\bm{x}}^o)]}{\sigma[R(\hat{\bm{x}})]}
= \frac{\mu_M - \E[R(\hat{\bm{x}}^o)]}{\sigma_M}
\end{equation}
である。よって，
\begin{align}
\mu_A &= \mu_M + \frac{\mu_M - \E[R(\hat{\bm{x}}^o)]}{\sigma_M}
\cdot\frac{\sigma_{AM}-\sigma_M^2}{\sigma_M}
= \E[R(\hat{\bm{x}}^o)]
+ \frac{\mu_M - \E[R(\hat{\bm{x}}^o)]}{\sigma_M^2}\sigma_{AM}\notag\\
&= \E[R(\hat{\bm{x}}^o)] + \frac{\sigma_{AM}}{\sigma_M^2}
\bigl(\mu_M - \E[R(\hat{\bm{x}}^o)]\bigr)
= \E[R(\hat{\bm{x}}^o)] + \frac{\sigma_A\rho_{AM}}{\sigma_M}
\bigl(\mu_M - \E[R(\hat{\bm{x}}^o)]\bigr)
\label{eq:zero-beta-capm}
\end{align}
が得られる。\eqref{eq:zero-beta-capm}は，CAPMにおける証券市場線を示す式
\begin{equation}
\mu_A = r + \frac{\sigma_{AM}}{\sigma_M^2}(\mu_M - r)
= r + \frac{\sigma_A\rho_{AM}}{\sigma_M}(\mu_M - r)
\end{equation}
における無リスク金利 $r$ を $\E[R(\hat{\bm{x}}^o)]$ に置き換えた式に相当し，
CAPMの $\beta$ に対応するのは
$\sigma_{AM}/\sigma_M^2 = \sigma_A\rho_{AM}/\sigma_M$ である。これらの議論を
\textbf{ゼロベータCAPM}という。

\begin{remark}
ゼロベータCAPMは定理\ref{thm:zero-beta-tangent}に基づくので，リスク性資産のみ
からなる任意の最小分散ポートフォリオ $\hat{\bm{x}}$(ただし，大域的最小分散
ポートフォリオ $\hat{\bm{x}}^*$ を除く)に対して適用できて，CAPMの無リスク金利
$r$ に相当する $\E[R(\hat{\bm{x}}^o)]$ も $\hat{\bm{x}}$ に応じて異なる。
\end{remark}

\begin{remark}
もしもCAPMと対応させるべく $\hat{\bm{x}}$ として市場ポートフォリオを選択すれば，
それに対応する $\E[R(\hat{\bm{x}}^o)]$ が無リスク金利 $r$ の役割を果たすことに
なる。その意味で，\eqref{eq:zero-beta-capm}は無リスク資産が存在しない場合の
CAPMの拡張になっている。
\end{remark}

\begin{remark}
しかし，定理\ref{thm:zero-beta-tangent}の接線が効率的フロンティアになるわけでは
なく，また，その接線上のポートフォリオが効率的ポートフォリオになるわけでもない。
\end{remark}

% ---------------- 演習問題 ----------------
\section{演習問題}

\begin{exercise}[問題1]
分散投資効果のケース1の設定(すべての投資先の収益率は同じ分布に従い互いに独立，
投資比率はすべて $1/N$，ポートフォリオの収益率は加重平均，リスクは標準偏差)で，
設定1と2を以下のように変更する。
\begin{itemize}
\item[1''.] 投資先の収益率 $R_j$， $j=1,\dots,N$ は次の式で与えられる：
\[
R_j = \beta_j R + \varepsilon_j.
\]
ただし，$R$ と $\varepsilon_j$ は互いに独立な確率変数で，$\varepsilon_j$ と
$\varepsilon_k$ も互いに独立とする。また
$\E[R]=\mu$， $\Var[R]=\sigma^2$， $\E[\varepsilon_j]=0$，
$\Var[\varepsilon_j]=\sigma_\varepsilon^2$ とする。
\end{itemize}
このとき，ポートフォリオの収益率の期待値と標準偏差を求めよ。
\end{exercise}

\begin{answerbox}
\textbf{考え方。}ケース2との違いは，共通因子 $R$ への感応度 $\beta_j$ が
\emph{銘柄ごとに異なる}ことである。計算の構造はケース2と同じで，共通成分と
固有成分に分けて期待値・分散を計算する。

\medskip
\textbf{解答。}ポートフォリオの収益率は
\[
R_p = \sum_{j=1}^{N}\frac{1}{N}R_j
= \frac{1}{N}\sum_{j=1}^{N}\beta_j R + \frac{1}{N}\sum_{j=1}^{N}\varepsilon_j.
\]
期待値は(期待値の線形性より，これは常に成り立つ等号)
\[
\E[R_p] = \frac{1}{N}\sum_{j=1}^{N}\beta_j\,\E[R]
+ \frac{1}{N}\sum_{j=1}^{N}\E[\varepsilon_j]
= \Bigl(\frac{1}{N}\sum_{j=1}^{N}\beta_j\Bigr)\mu.
\]
分散は，$R$ と $\varepsilon_j$ たちがすべて独立(無相関)であることから
\[
\Var[R_p] = \Bigl(\frac{1}{N}\sum_{j=1}^{N}\beta_j\Bigr)^2\sigma^2
+ \frac{1}{N^2}\sum_{j=1}^{N}\sigma_\varepsilon^2
= \Bigl(\frac{1}{N}\sum_{j=1}^{N}\beta_j\Bigr)^2\sigma^2
+ \frac{\sigma_\varepsilon^2}{N}
\]
となる(「和の分散 $=$ 分散の和」は常に成り立つ等号ではなく，無相関だから
成り立つ)。よって標準偏差は
\[
\sigma_p = \sqrt{\bar\beta^2\sigma^2 + \frac{\sigma_\varepsilon^2}{N}},
\qquad
\bar\beta \equiv \frac{1}{N}\sum_{j=1}^{N}\beta_j.
\]
\end{answerbox}

\begin{exercise}[問題2]
問題1で，$N$ を無限に大きくしていくと，リスクはある値に近付いていく(モデルで
与えている平均，分散，係数などはすべて既知とする)。このとき，極めて多様な資産が
存在して，それらを極めて上手に組み合わせれば，$N$ を無限大にした極限で，全体の
リスクをゼロにすることができる。そのような理想的な極限を考えるとき，(リスクで
なく)期待収益率はどうなるか。
\end{exercise}

\begin{answerbox}
\textbf{考え方。}問題1の結果で $N\to\infty$ の極限をとる。リスクの極限値は
$\bar\beta$(ベータの平均)に依存するので，「リスクをゼロにする組み合わせ」とは
$\bar\beta\to 0$ とすることに他ならない。そのとき期待収益率がどうなるかを見る。

\medskip
\textbf{解答。}問題1の結果より，
\[
\sigma_p = \sqrt{\bar\beta^2\sigma^2 + \frac{\sigma_\varepsilon^2}{N}}
\;\xrightarrow{N\to\infty}\; |\bar\beta|\,\sigma
\qquad\Bigl(\text{ただし }
\bar\beta = \frac{1}{N}\sum_{j=1}^N\beta_j \xrightarrow{N\to\infty}
\beta < \infty \text{ とする}\Bigr).
\]
この結果は，$\beta_j$ の組み合わせ次第でリスクを低減できることを示唆している。
例えば，$\beta_j$ の符号が異なる資産を非常にうまく組み合わせれば，$\bar\beta$ を
0に近付けることもできる。

では，その場合，期待収益率はどうなるか。$\bar\beta$ を0に漸近させれば全体の
リスクを0にできるが，そのとき
\[
\E[R_p] = \bar\beta\,\mu \;\longrightarrow\; 0
\]
より，\textbf{期待収益率もゼロになる}。

この場合，期待収益率が0で，リスク($=$変動性)も0なので，本来は確率変数である
はずの収益率自体が0になっていく。すなわち，
\textbf{リスクをとらなければリターンも得られない}ことを示唆している。
\end{answerbox}
